\section{Reduction and scan}
\label{sec:scan}

Reduction and scan are the two families where the roof is entirely a bandwidth
question, so the ladder is a ladder of passes over the input rather than of
arithmetic. Like the FFT family, the mechanisms are implemented and tested and the
throughput rows are pending.

\subsection{The ladder is about passes, not instructions}

The scan ladder has five rungs. The first is a Hillis-Steele block scan, which is
work inefficient by construction: $O(N \log N)$ work at block scope. The second is
Blelloch's work efficient upsweep and downsweep with bank conflict
padding~\cite{blelloch1990}, which does $O(N)$ work and pays for it in barriers.
Those two are block primitives, not device wide scans, so both are wrapped in the
same scan-then-propagate decomposition the third rung uses; wrapping them in
different ones would have made the difference between the rows two changes rather
than one. The third rung moves about $4N$ bytes, the fourth (reduce then scan) about
$3N$, and the fifth, decoupled look-back, about $2N$, which is the lower bound for
the problem.

Decoupled look-back~\cite{merrill2016scan} is the top rung and it carries three
constraints that are not optional. The tile index comes from an atomic increment on
a global counter and never from the block index, because a block that claimed index
$k$ did so after every block holding a lower index had already started, so those
blocks are resident and can make progress; indexing by block index would let a
resident block wait on one the scheduler has not launched, which deadlocks. The
status word is published with release semantics and read with acquire semantics,
with the aggregate and the inclusive prefix sitting beside the flag as plain
stores, which is exactly what release and acquire exist to order. And the status
array is zeroed by a memset enqueued ahead of the kernel and never inside it,
because a block cannot zero its own slot when a predecessor may already have read
it. That memset is real traffic and the plan counts it in the model rather than
leaving it out to make the top rung look cheaper.

\subsection{Determinism has a price and it is measured, not defended}

Floating point addition is not associative, so a reduction gives a different answer
when it combines its partials in a different order, and every fast rung above lets
the launch decide that order. That is not a defect; it is what a fast reduction
costs.

A deterministic rung exists for both families and makes the answer bit identical
three ways. The block count is a pure function of the length and never asks the
occupancy API, so the same input gives the same grid on a busy device, an idle one,
or a different card. The element to thread mapping follows from that grid, and each
thread combines its own elements in index order. And the two levels are combined by
an explicit halving tree, so the shape of the tree is a property of the block size
rather than of the schedule. Decoupled look-back cannot make that promise, because
how many predecessors a tile has to combine depends on which of them happened to
finish first, which is why the deterministic scan is reduce then scan and is a
separate name rather than an alias.

The artifact for that claim is committed:
\texttt{experiments/results/scan\_determinism.txt} holds the output hashes from ten
consecutive runs at $N = 2^{24}$ and $N = 2^{28}$. It is a determinism artifact
produced at unlocked clocks, not a performance artifact, and it is valid as such:
the value it records is a property of the arithmetic and has nothing to do with how
fast the machine was running. The cost of determinism is a benchmark row of its own
next to the fast rungs, and that row is pending.

\begin{table}[htbp]
\centering
\begin{tabular}{lllrr}
\toprule
rung & mechanism & bytes moved & effective GB/s at the largest swept size & percent of CUB \\
\midrule
s0 Hillis-Steele & block scan, work inefficient & N log N work at block scope & 509.9 & 44.3 \\
s1 Blelloch & upsweep and downsweep, bank conflict padding & N work, twice the barriers & 436.7 & 37.9 \\
s2 three kernel & scan then propagate & about 4 N bytes & 576.9 & 50.1 \\
s3 reduce then scan & two passes over the input & about 3 N bytes & 595.2 & 68.9 \\
s4 decoupled look-back & single pass, device scope handshake & about 2 N bytes & 551.7 & 95.5 \\
deterministic scan & fixed grid, index order, explicit halving tree & about 3 N bytes & 595.2 & 68.9 \\
\bottomrule
\end{tabular}

\caption{The reduction and scan ladders. Bytes moved per element is derived from
the decomposition; throughput and percent of CUB are pending the sweep.}
\label{tab:scan}
\end{table}
